\chapter*{부록 D. AI 거버넌스 성숙도 평가 도구 상세 지표}
\label{ch:appendix_d}
\addcontentsline{toc}{chapter}{부록 D. AI 거버넌스 성숙도 평가 도구 상세 지표}
\markboth{부록 D}{부록 D}

\section*{D.0 FAIR-D 프레임워크: AI 고유 리스크 성숙도 모델}

기존의 AI 거버넌스 성숙도 모델들은 대부분 CMMI나 DCAM과 같은 전통적 IT 성숙도 모델을 AI에 그대로 이식한 것이다. 이 접근의 한계는 명확하다: AI는 전통적 소프트웨어와 근본적으로 다른 위험을 내포한다. 코드는 배포 후 변하지 않지만, AI 모델은 \textbf{소리 없이 드리프트}한다. 일반 소프트웨어는 결정 근거를 추적할 수 있지만, AI 모델은 \textbf{설명할 수 없는 결과}를 생성한다. 전통적 시스템은 없는 데이터를 만들어내지 않지만, 생성형 AI는 \textbf{그럴듯한 허구(할루시네이션)}를 생성한다.

본 부록은 이러한 인식에서 출발한 \textbf{FAIR-D 프레임워크(Fairness · Autonomy · Interpretability · Reliability · Drift)}를 제시한다. AI 거버넌스에서만 의미 있는 5개의 핵심 벡터를 중심으로 조직의 성숙도를 진단하는 이 모델은 NIST AI RMF(2023), ISO/IEC 42001(2023), 그리고 2026년 시행된 한국 AI 기본법의 고영향 AI 요구사항을 통합하여 설계되었다.

\begin{table}[htbp]
\centering
\caption{FAIR-D 프레임워크: 5대 핵심 벡터}
\label{tab:faird_vectors}
\begin{tabularx}{\textwidth}{|l|l|X|}
\hline
\textbf{벡터} & \textbf{의미} & \textbf{AI 고유성} \\ \hline
\textbf{F}airness & 공정성 & 알고리즘적 편향은 전통 SW에 없는 AI 고유 위험이다 \\ \hline
\textbf{A}utonomy Control & 자율성 통제 & AI의 자율적 의사결정 범위를 인간이 실질적으로 통제하는가 \\ \hline
\textbf{I}nterpretability & 해석가능성 & 모델의 결정을 인간이 이해하고 설명할 수 있는가 \\ \hline
\textbf{R}eliability & 신뢰성·할루시네이션 통제 & 생성형 AI의 허구 생성을 감지하고 방지하는가 \\ \hline
\textbf{D}rift Management & 드리프트 관리 & 배포 후 성능 저하를 자동으로 감지하고 대응하는가 \\ \hline
\end{tabularx}
\end{table}

각 벡터는 1단계(부재)부터 5단계(선도)까지 독립적으로 평가되며, 5개 벡터의 프로파일이 조직의 ``AI 거버넌스 DNA''를 형성한다. 5개 벡터의 단순 합산이 아니라 \textbf{프로파일 패턴}이 중요하다. 예를 들어 공정성은 5단계이지만 드리프트 관리가 1단계인 조직은, 높은 평균 점수에도 불구하고 심각한 운영 리스크를 안고 있다.

\section*{D.1 벡터별 성숙도 세부 지표}

\subsection*{D.1.1 Fairness (공정성 거버넌스)}

\begin{table}[htbp]
\centering
\caption{공정성 거버넌스 성숙도 5단계}
\label{tab:fairness_maturity}
\begin{tabularx}{\textwidth}{|l|l|X|}
\hline
\textbf{단계} & \textbf{명칭} & \textbf{특성 및 진단 기준} \\ \hline
1단계 & 부재(Absent) & 공정성 개념이 없음. 데이터 품질만 고려. 편향 테스트 없음 \\ \hline
2단계 & 인식(Aware) & 공정성 중요성 인식. 비공식 검토. 측정 기준 미정립 \\ \hline
3단계 & 측정(Measured) & 인구통계 공정성(DP), 균등 기회(EO) 등 정량 지표 적용. 정기 평가 수행 \\ \hline
4단계 & 통합(Integrated) & CI/CD 파이프라인에 공정성 게이트 내장. 자동 회귀 탐지. 재학습 트리거 \\ \hline
5단계 & 선도(Leading) & 교차집단(Intersectional) 공정성 분석. 알고리즘 공정성 감사 결과 외부 공개. 사회적 피드백 루프 \\ \hline
\end{tabularx}
\end{table}

\textbf{핵심 진단 문항}:
\begin{enumerate}
    \item 배포된 모든 AI 시스템에 대해 인구통계적 공정성(Demographic Parity) 지표가 정기적으로 측정되는가?
    \item 공정성 지표가 허용 임계값(예: DP difference $<$ 0.1)을 초과할 경우 자동으로 알람이 발송되는가?
    \item 데이터 수집 단계부터 특정 집단의 과소 대표(Underrepresentation) 여부를 검토하는 프로세스가 있는가?
    \item 공정성 테스트 결과가 배포 승인 게이트의 필수 통과 조건인가?
    \item 공정성 감사 결과가 외부에 공개되거나 규제 기관에 보고되는가?
\end{enumerate}

\subsection*{D.1.2 Autonomy Control (자율성 통제)}

AI 시스템이 인간의 실질적 감독 없이 의사결정을 내리는 범위가 확대될수록, 자율성 통제의 성숙도가 중요해진다. 이는 특히 한국 AI 기본법의 ``고영향 AI에 대한 인간 감독 의무''와 직결된다.

\begin{table}[htbp]
\centering
\caption{자율성 통제 성숙도 5단계}
\label{tab:autonomy_maturity}
\begin{tabularx}{\textwidth}{|l|l|X|}
\hline
\textbf{단계} & \textbf{명칭} & \textbf{특성 및 진단 기준} \\ \hline
1단계 & 방임(Unchecked) & AI 결정이 인간 검토 없이 자동 집행됨. 무제한 자동화 \\ \hline
2단계 & 사후검토(Post-hoc) & 자동 집행 후 인간이 사후 감사. 이미 발생한 오류 교정 불가 \\ \hline
3단계 & 게이트형(Gated) & 고위험 결정은 인간이 사전 승인. 위험 수준 분류 체계 존재 \\ \hline
4단계 & 동적 감독(Dynamic) & 신뢰도 기반 자동 에스컬레이션. 불확실성 높을 때 자동으로 인간 개입 \\ \hline
5단계 & 설계 내재화(By-design) & Human-in-the-loop가 아키텍처에 내재화. 인간 감독 우회 불가 구조 \\ \hline
\end{tabularx}
\end{table}

\textbf{핵심 진단 문항}:
\begin{enumerate}
    \item 조직 내 AI 시스템의 의사결정을 ``완전 자동화'', ``인간 검토 필요'', ``인간 최종 결정'' 세 등급으로 분류한 목록이 있는가?
    \item 고영향 AI(채용, 대출, 의료 진단 등)에서 AI 권고를 인간이 거부할 수 있는 실질적 절차가 보장되는가?
    \item AI 시스템이 인간의 개입 없이 내릴 수 있는 결정의 상한(금액, 영향 범위 등)이 명시적으로 정의되어 있는가?
    \item 자동화 시스템이 예외적 상황을 감지하면 자동으로 인간에게 에스컬레이션하는 메커니즘이 있는가?
    \item AI 기본법상 고영향 AI에 해당하는 시스템 목록과 각각의 인간 감독 절차가 문서화되어 있는가?
\end{enumerate}

\subsection*{D.1.3 Interpretability (해석가능성·설명가능성)}

\begin{table}[htbp]
\centering
\caption{해석가능성 성숙도 5단계}
\label{tab:interpretability_maturity}
\begin{tabularx}{\textwidth}{|l|l|X|}
\hline
\textbf{단계} & \textbf{명칭} & \textbf{특성 및 진단 기준} \\ \hline
1단계 & 블랙박스(Opaque) & 설명 없이 결과만 제공. XAI 도구 없음 \\ \hline
2단계 & 사후 설명(Post-hoc) & SHAP/LIME 등 사후 해석 도구 보유. 비체계적 적용 \\ \hline
3단계 & 표준화(Standardized) & 모델 카드 작성 의무화. 주요 변수 중요도 정기 보고 \\ \hline
4단계 & 대화형(Interactive) & 의사결정자가 특정 케이스의 결정 근거를 실시간 조회 가능 \\ \hline
5단계 & 감사가능(Auditable) & 전체 추론 체인 추적 가능. 규제 기관 요청 시 즉시 제출 가능 \\ \hline
\end{tabularx}
\end{table}

\textbf{핵심 진단 문항}:
\begin{enumerate}
    \item 배포된 모든 모델에 대해 모델 카드(Model Card)가 작성·관리되는가?
    \item 영향을 받은 개인이 AI 결정의 주요 근거를 이해할 수 있도록 설명을 제공하는가?
    \item XAI 도구(SHAP, LIME 등)가 개발 파이프라인에 통합되어 정기적으로 실행되는가?
    \item 고위험 결정에서 AI가 참조한 주요 특성(Feature)과 그 기여도를 감사 로그에 기록하는가?
    \item 규제 기관이 특정 결정의 근거를 요청할 경우 48시간 내에 제출할 수 있는 체계가 있는가?
\end{enumerate}

\subsection*{D.1.4 Reliability / 할루시네이션 통제 (생성형 AI 신뢰성)}

이 벡터는 전통적 AI 성숙도 모델에 없던 차원이다. 생성형 AI가 조직 업무에 광범위하게 도입되면서, 할루시네이션(사실과 다른 그럴듯한 출력) 통제는 독립적 거버넌스 영역이 되었다.

\begin{table}[htbp]
\centering
\caption{할루시네이션 통제 성숙도 5단계}
\label{tab:reliability_maturity}
\begin{tabularx}{\textwidth}{|l|l|X|}
\hline
\textbf{단계} & \textbf{명칭} & \textbf{특성 및 진단 기준} \\ \hline
1단계 & 무방비(Unguarded) & 생성 결과를 그대로 신뢰. 사실성 검증 없음 \\ \hline
2단계 & 경고형(Warned) & 사용자에게 AI 출력 검증 안내. 체계적 검증 방법 미비 \\ \hline
3단계 & 검증형(Verified) & RAG 기반 출처 인용. 핵심 정보에 대한 사실 검증 파이프라인 \\ \hline
4단계 & 자동화형(Automated) & 실시간 팩트체크 도구 통합. 할루시네이션 탐지 점수화 및 임계값 관리 \\ \hline
5단계 & 신뢰보증(Assured) & 출처 가중치 관리, 불확실성 정량화, 자신감 점수 제공. 규제 기관 증명 가능 \\ \hline
\end{tabularx}
\end{table}

\textbf{핵심 진단 문항}:
\begin{enumerate}
    \item 생성형 AI 시스템이 제공하는 정보에 대한 출처 인용(Citation) 또는 근거 문서 링크가 제공되는가?
    \item 할루시네이션 탐지 비율이 정기적으로 측정되고 목표 임계값이 설정되어 있는가?
    \item 생성형 AI 출력이 고위험 결정(의료 조언, 법적 판단, 금융 추천)에 직접 사용될 경우 추가 검증 단계가 의무화되어 있는가?
    \item 사용자가 AI 출력의 신뢰도 수준을 확인할 수 있는 인터페이스가 제공되는가?
    \item 할루시네이션 인시던트 발생 시 원인 분석 및 재발 방지 절차가 있는가?
\end{enumerate}

\subsection*{D.1.5 Drift Management (드리프트 관리)}

모델 드리프트는 AI 특유의 ``조용한 실패''다. 전통적 소프트웨어는 배포 후 동일한 입력에 동일한 출력을 낸다. AI 모델은 현실 데이터 분포가 변화하면 학습 당시와 다른 성능을 보이며, 이는 심각한 경우 사회적 피해로 이어진다.

\begin{table}[htbp]
\centering
\caption{드리프트 관리 성숙도 5단계}
\label{tab:drift_maturity}
\begin{tabularx}{\textwidth}{|l|l|X|}
\hline
\textbf{단계} & \textbf{명칭} & \textbf{특성 및 진단 기준} \\ \hline
1단계 & 미인식(Unmonitored) & 드리프트 개념 없음. 초기 배포 성능이 영구적이라 가정 \\ \hline
2단계 & 반응적(Reactive) & 성능 저하 보고 후 사후 조사. 임계값 없음 \\ \hline
3단계 & 감지형(Detected) & 정기적 성능 지표 모니터링. PSI/KS 통계 기반 드리프트 탐지 \\ \hline
4단계 & 자동화형(Automated) & 드리프트 탐지 시 자동 알람 및 재학습 트리거. 공정성 드리프트 별도 추적 \\ \hline
5단계 & 예측형(Predictive) & 데이터 변화 조기 신호 감지. 재학습 전 성능 영향 시뮬레이션 가능 \\ \hline
\end{tabularx}
\end{table}

\textbf{핵심 진단 문항}:
\begin{enumerate}
    \item 배포된 모든 모델에 대해 성능 지표(정확도, F1 등)가 실시간 또는 일별로 모니터링되는가?
    \item 데이터 드리프트(PSI, KS 통계 등)와 모델 드리프트(성능 저하)를 독립적으로 추적하는가?
    \item 공정성 지표의 드리프트(특정 집단에 대한 성능 저하)를 전체 성능 드리프트와 별도로 모니터링하는가?
    \item 드리프트 임계값 초과 시 자동으로 재학습이 트리거되거나 담당자에게 알람이 발송되는가?
    \item 마지막 재학습 이후 경과 시간과 데이터 변화량이 자동으로 추적·기록되는가?
\end{enumerate}

\section*{D.2 FAIR-D 성숙도 진단 시트}

각 벡터별로 1\textasciitilde{}5점을 부여하고, 5각형 레이더 차트로 시각화하면 조직의 AI 거버넌스 ``DNA 프로파일''을 확인할 수 있다.

\begin{table}[htbp]
\centering
\caption{FAIR-D 진단 결과 기록 시트}
\label{tab:faird_scoresheet}
\begin{tabularx}{\textwidth}{|l|l|l|X|}
\hline
\textbf{벡터} & \textbf{현재 단계} & \textbf{목표 단계} & \textbf{핵심 개선 과제} \\ \hline
F: 공정성 & \hspace{1.5cm} / 5 & \hspace{1.5cm} / 5 & \\ \hline
A: 자율성 통제 & \hspace{1.5cm} / 5 & \hspace{1.5cm} / 5 & \\ \hline
I: 해석가능성 & \hspace{1.5cm} / 5 & \hspace{1.5cm} / 5 & \\ \hline
R: 신뢰성/할루시네이션 & \hspace{1.5cm} / 5 & \hspace{1.5cm} / 5 & \\ \hline
D: 드리프트 관리 & \hspace{1.5cm} / 5 & \hspace{1.5cm} / 5 & \\ \hline
\textbf{합계} & \hspace{1.5cm} / 25 & \hspace{1.5cm} / 25 & \\ \hline
\end{tabularx}
\end{table}

\section*{D.3 산업별 FAIR-D 목표 프로파일 (2026년 권고)}

산업 특성에 따라 각 벡터의 중요도가 다르다. 아래 표는 주요 산업에서 권고되는 최소 성숙도 수준이며, 한국 AI 기본법의 고영향 AI 분류와 EU AI Act의 고위험 AI 분류를 반영하였다.

\begin{table}[htbp]
\centering
\caption{산업별 FAIR-D 최소 권고 성숙도 (1\textasciitilde{}5단계)}
\label{tab:faird_industry}
\begin{tabularx}{\textwidth}{|l|l|l|l|l|l|}
\hline
\textbf{산업} & \textbf{F} & \textbf{A} & \textbf{I} & \textbf{R} & \textbf{D} \\ \hline
의료·헬스케어 & 4 & \textbf{5} & \textbf{5} & 4 & 4 \\ \hline
금융·신용평가 & \textbf{5} & 4 & 4 & 3 & \textbf{5} \\ \hline
공공행정·복지 & \textbf{5} & \textbf{5} & 4 & 3 & 4 \\ \hline
채용·인사 & \textbf{5} & 4 & \textbf{5} & 3 & 3 \\ \hline
생성형 AI 서비스 & 4 & 3 & 3 & \textbf{5} & 4 \\ \hline
제조·산업안전 & 3 & \textbf{5} & 4 & 2 & \textbf{5} \\ \hline
마케팅·추천 & 4 & 2 & 3 & 3 & 4 \\ \hline
\end{tabularx}
\end{table}

\textbf{굵게 표시된 항목}: 해당 산업에서 AI 기본법 또는 EU AI Act의 의무 요구 수준.

\section*{D.4 기존 성숙도 모델과 FAIR-D의 비교}

\begin{table}[htbp]
\centering
\caption{기존 AI 거버넌스 성숙도 모델과 FAIR-D 비교}
\label{tab:faird_comparison}
\begin{tabularx}{\textwidth}{|l|X|X|}
\hline
\textbf{구분} & \textbf{기존 모델 (CMMI 이식형)} & \textbf{FAIR-D (AI 고유형)} \\ \hline
설계 철학 & 일반 IT 프로세스 성숙도를 AI에 적용 & AI 고유 위험에서 역설계 \\ \hline
평가 단위 & 전사 거버넌스 전반 & AI 고유 리스크 벡터별 독립 진단 \\ \hline
할루시네이션 & 미포함 & 독립 벡터(R)로 측정 \\ \hline
드리프트 & 부분 포함 (모니터링 항목) & 독립 벡터(D)로 세분화 \\ \hline
공정성 드리프트 | 미분리 & 드리프트(D)와 공정성(F) 교차 추적 \\ \hline
결과 활용 & 단일 점수 또는 레벨 & 5차원 프로파일 + 산업별 벤치마크 \\ \hline
규제 연계 & 일반적 컴플라이언스 & AI 기본법 고영향 AI + EU AI Act 직접 매핑 \\ \hline
\end{tabularx}
\end{table}

\begin{tcolorbox}[colback=blue!5!white, colframe=blue!75!black, title={\textbf{FAIR-D 활용 가이드}}]
\begin{enumerate}
    \item \textbf{진단}: 5개 벡터를 각각 독립적으로 평가한다 (1\textasciitilde{}5점).
    \item \textbf{프로파일링}: 5각형 레이더 차트로 시각화하여 취약 벡터를 식별한다.
    \item \textbf{산업 벤치마크}: 자사 산업의 권고 최소 수준(D.3 표)과 현재 수준을 비교한다.
    \item \textbf{우선순위 결정}: 가장 낮은 벡터 또는 산업 필수 벡터에서 가장 큰 갭이 있는 부분부터 개선한다.
    \item \textbf{분기 재평가}: 3개월마다 재진단하여 변화 추이를 추적한다.
\end{enumerate}
FAIR-D는 점수가 아니라 \textbf{대화를 위한 도구}다. 각 벡터의 점수보다, 왜 그 단계에 머물러 있는지에 대한 조직 내 토론이 더 중요하다.
\end{tcolorbox}
